% TOP_PROBLEM: {"id": 1157, "title": "Pillai's Conjecture", "category_id": 1, "category": {"id": 1, "name": "number_theory", "display_name": "Number Theory", "description": "Properties of integers, prime numbers, Diophantine equations.", "slug": "number-theory", "order_index": 1, "created_at": "2026-07-31T15:26:25.670Z"}, "status": "open"}
% FIRST_POSTED: null
% ENTRY_KIND: statement_only
% Imported from: batch04.tex; UnsolvedMath ID 1157
\documentclass[11pt]{article}
\usepackage[margin=25mm]{geometry}
\usepackage{fontspec}
\setmainfont{Latin Modern Roman}
\usepackage{amsmath,amssymb,mathtools}
\usepackage{unicode-math}
\setmathfont{Latin Modern Math}
\usepackage{xurl,hyperref}
\hypersetup{colorlinks=true,urlcolor=blue}
\setcounter{secnumdepth}{0}
\setlength{\parindent}{0pt}
\setlength{\parskip}{5pt}
\setlength{\emergencystretch}{3em}
\newcommand{\sref}[2]{\hyperlink{source:#1:#2}{[S#2]}}
\newcommand{\cataloguescope}[1]{\paragraph{Recorded target.}#1}
\begin{document}
% UnsolvedMath ID 1157; source catalogue code NT-023
\setcounter{section}{1156}
\section{Pillai's Conjecture}\label{1157}
{\small\sffamily UnsolvedMath ID 1157 · Number Theory}
\subsection{Definitions and mathematical statement}

\paragraph{Shared notation.}$\mathbb N=\{1,2,\ldots\}$, and $\mathbb Z,\mathbb Q,\mathbb R,\mathbb C$ denote the integers, rationals, reals and complex numbers. Any other conventions are specified locally.


For a fixed positive integer \(k\), let
\[
 S_k=\{(m,n)\in\mathbb N^2:|2^m-3^n|=k\}.
\]
Both exponents are positive, and both signs of the difference are allowed.
The bases two and three stay fixed as \(k\) varies.
\paragraph{Statement recorded in the supplied source.}
For every \(k\in\mathbb N\), is \(S_k\) finite?
The fixed-base assertion is a consequence of the finite-rank
multiplicative-group equation theorem cited below: for each sign,
\(2^m-3^n=\pm k\) becomes a two-variable nondegenerate linear equation
with variables \(2^m,3^n\).  This question does not allow the bases
themselves to vary. \sref{1157}{2}
\subsection{Short English statement}
For each fixed positive difference, can only finitely many powers of two and powers of three have that difference in absolute value?
\subsection{Sources}
\begin{description}
\item[\hypertarget{source:1157:1}{S1}] ULAM.AI and the UnsolvedMath Contributors, \emph{UnsolvedMath: A Curated Collection of Open Mathematics Problems}, record 1157 (NT-023). CC BY 4.0. \url{https://huggingface.co/datasets/ulamai/UnsolvedMath}.
\item[\hypertarget{source:1157:2}{S2}] Jan-Hendrik Evertse, Hans Peter Schlickewei and Wolfgang M. Schmidt, \emph{Linear equations in variables which lie in a multiplicative group}, Annals of Mathematics 155 (2002), 807--836, Theorem 1.1; arXiv:math/0409604. \url{https://arxiv.org/abs/math/0409604}.
\end{description}
\label{end:1157}
\end{document}
